% Chapter 4, Section 4.1 -- common protocol only; no performance results.
\subsection{Experimental protocol}
\label{subsec:experimental_setup}

All experiments followed the battery-level split, trajectory-concentrated label-allocation protocol, and validation-based model-selection procedure defined in Sections~\ref{sec:data} and~\ref{sec:method}. All estimators used the same 14 charging features, 40-cycle input window, training/validation/test split, label-allocation mask, target scaling, and fully held-out test cells. The evaluated SOH-label budgets were 100\%, 50\%, 30\%, 10\%, and 5\%, with 100\% serving as the full-supervision reference. For every sparse-label condition, the same allocation was paired across the compared estimators; the test cells remained unavailable during training, normalisation, early stopping, and hyperparameter selection. Deep-learning models were trained with comparable optimisation budgets and the same early-stopping principle.

Each configuration was evaluated across three paired repetitions, and all reported results are expressed as mean $\pm$ standard deviation. Point-estimation performance was quantified by the mean absolute error (MAE), root mean squared error (RMSE), and coefficient of determination ($R^2$). To assess trajectory quality beyond pointwise accuracy, the ablation and sensitivity analyses additionally report the monotonicity-violation rate, cumulative upward excess, and cycle-interval-normalised rate variation. All hyperparameters, including $\lambda_{\mathrm{mono}}$ and $\lambda_{\mathrm{rate}}$, were selected using the validation set only; the test set was used exclusively for final evaluation.
